% =========================================================
%  sesi/02/sesi.tex
%  Sesi 2 — Siklus permainan dan pengendali waktu
% =========================================================

\sesi{02}{Siklus Permainan dan Pengendali Waktu}{1.1}{110 menit}

Sesi pertama berakhir pada jendela yang diam. Latar diisi satu kali, lalu
program menunggu peristiwa dengan \code{pygame.event.wait()}; selama pemain
tidak menyentuh apa pun, layar tidak berubah dan tidak ada pekerjaan yang
dilakukan. Jendela semacam itu memadai untuk menampilkan informasi, tetapi
belum memadai untuk permainan.

Pada minggu 2 ini jendela tersebut mulai hidup. Anda menulis \istilah{game
loop}: perulangan yang membaca masukan, memperbarui keadaan, dan menggambar
ulang jendela secara berulang. Agar pengulangan itu tidak memakan seluruh
kemampuan prosesor, Anda menambahkan pengendali waktu dan menampilkan laju
bingkai (\istilah{FPS}) pada judul jendela. Angka itu menjadi alat ukur pertama
Anda untuk menilai apakah program bekerja sebagaimana mestinya.

% ---------------------------------------------------------
\bagian{Tujuan Sesi}

\begin{tujuanbox}
Setelah menuntaskan sesi ini, Anda dapat menunjukkan:
\begin{enumerate}[leftmargin=1.5em, itemsep=0.2em]
  \item jendela \NamaPermainan\ yang berjalan terus dan judulnya memuat angka
        FPS yang berubah dari waktu ke waktu;
  \item jendela yang tetap menanggapi masukan ketika pemain mendiamkan tombol
        dan tetikus, tanpa tanda ``tidak menanggapi'';
  \item jendela yang berhenti dengan rapi ketika ditutup lewat tombol silang
        atau tombol ESC, dan terminal kembali menampilkan
        \code{Program selesai.};
  \item perbedaan angka pada judul jendela ketika program dijalankan dengan
        batas 60 FPS dan 30 FPS;
  \item penjelasan, sambil menunjuk angka pada judul jendela, mengapa laju
        yang diminta tidak selalu sama dengan laju yang benar-benar tercapai.
\end{enumerate}
\end{tujuanbox}

% ---------------------------------------------------------
\bagian{Prasyarat}

\begin{prasyaratbox}
\begin{itemize}[leftmargin=1.4em, itemsep=0.25em]
  \item Checkpoint sesi 1 pada \file{code/sesi-01} sudah berjalan: jendela
        terbuka, latar terisi, dan tombol silang menutup program dengan rapi.
  \item Anda telah membaca Bab IV buku ajar, khususnya bagian kerangka siklus,
        pembacaan peristiwa, dan pengendali waktu.
  \item Bab II buku ajar cukup dibaca pada bagian yang membedakan langkah tetap
        dan langkah berubah; sesi ini baru memakainya sebagai konsep, belum
        sebagai perhitungan gerak.
  \item Python 3.8 atau lebih baru dan Pygame 2 masih terpasang pada penafsir
        yang sama dengan yang Anda pakai pada sesi 1.
\end{itemize}
\end{prasyaratbox}

% ---------------------------------------------------------
\bagian{Peta Kemampuan}

\petasesi
  {Jendela diam: latar diisi sekali, lalu program menunggu satu peristiwa
   dengan \code{event.wait()}.}
  {Siklus masukan--pembaruan--penayangan, \code{event.get()} dan
   \code{QUIT}, \code{Clock} pada 60 FPS, serta angka FPS pada judul jendela.}
  {Pemain berukuran 36 piksel yang bergerak dengan \code{Vector2} dan
   \istilah{delta time} pada sesi 3.}

% ---------------------------------------------------------
\bagian{Konsep yang Dipakai}

\subsection*{Tiga tahap yang diulang setiap bingkai}

Satu putaran \istilah{game loop} terdiri atas tiga tahap yang selalu muncul
dalam urutan yang sama.

\begin{enumerate}[leftmargin=1.6em, itemsep=0.2em]
  \item \textbf{Masukan.} Program membaca peristiwa yang masuk sejak putaran
        sebelumnya: jendela hendak ditutup, tombol ditekan, atau tetikus
        digeser.
  \item \textbf{Pembaruan.} Keadaan permainan diubah menurut masukan dan
        menurut waktu yang berjalan. Pada sesi ini belum ada objek yang
        dipindahkan, sehingga tahap ini masih kosong.
  \item \textbf{Penayangan.} Seluruh isi jendela digambar ulang menurut keadaan
        terbaru, lalu ditayangkan dengan \code{pygame.display.flip()}.
\end{enumerate}

\begin{center}
% figures/siklus-game-loop.tikz
% Siklus input -> update -> render dengan pengendali waktu.
\begin{tikzpicture}[
  tahap/.style={rectangle, rounded corners=2pt, draw=ubNavy, thick,
    fill=ubSoft, align=center, text width=2.5cm, minimum height=1.15cm,
    inner sep=3pt, font=\small},
  waktu/.style={tahap, draw=ubAccent, fill=white},
  panah/.style={-{Stealth[length=2.4mm]}, thick, draw=ubAccent}
]
  \node[tahap] (input) {\textbf{Input}\\ \footnotesize baca peristiwa};
  \node[tahap, right=1.0cm of input] (update) {\textbf{Update}\\ \footnotesize ubah keadaan};
  \node[tahap, right=1.0cm of update] (render) {\textbf{Render}\\ \footnotesize gambar ulang};
  \node[waktu, below=0.9cm of update] (clock) {\textbf{Clock}\\ \footnotesize ukur selang, batasi FPS};

  \draw[panah] (input) -- (update);
  \draw[panah] (update) -- (render);
  \draw[panah, dashed] (render.south) -- ++(0,-0.55) -| (clock.east);
  \draw[panah, dashed] (clock.west) -| (input.south);
  \draw[panah] (render.north) -- ++(0,0.45) -| (input.north);
  \node[font=\scriptsize\itshape, ubMuted, above=0.42cm of update]
    {satu bingkai, diulang terus};
\end{tikzpicture}

\end{center}

Masukan didahulukan karena pembaruan perlu mengetahui kehendak pemain pada
bingkai itu juga. Penayangan dikerjakan terakhir karena yang boleh dilihat
pemain hanyalah keadaan yang sudah selesai dihitung, bukan keadaan setengah
jadi.

\subsection*{Peristiwa dibaca, bukan ditunggu}

Pygame menaruh setiap kejadian ke dalam antrean. Perintah
\code{pygame.event.get()} mengambil seluruh isi antrean tanpa menunggu kejadian
baru, sehingga siklus tetap berjalan walaupun pemain sedang diam. Sebaliknya,
\code{pygame.event.wait()} menghentikan program sampai ada peristiwa, dan
karena itu tidak boleh dipakai di dalam siklus. Setiap peristiwa memiliki jenis;
jenis \code{QUIT} berarti pemain meminta jendela ditutup.

\subsection*{Anggaran waktu satu bingkai}

Batas laju 60 FPS berarti program ingin menyelesaikan paling banyak 60 bingkai
dalam satu detik. Jatah waktu untuk satu bingkai karena itu adalah

\[
\Delta t_{\text{jatah}} = \frac{1000\ \text{milidetik}}{60}
\approx 16{,}7\ \text{milidetik}.
\]

Lambang $\Delta t_{\text{jatah}}$ menyatakan jatah waktu satu bingkai, angka
1000 menyatakan jumlah milidetik dalam satu detik, dan angka 60 menyatakan
batas bingkai per detik yang diminta. Apabila seluruh pekerjaan satu bingkai
selesai dalam 16,7 milidetik atau kurang, \code{clock.tick()} menahan sisa
waktunya. Apabila pekerjaan itu memerlukan 30 milidetik, jatahnya sudah
terlampaui: \code{tick()} tidak menahan apa pun dan laju yang tercapai hanya
sekitar 33 FPS, bukan 60. Inilah sebabnya angka pada judul jendela perlu
dibaca sebagai kenyataan, bukan sebagai janji.

\subsection*{Langkah tetap dan langkah berubah}

Pada \istilah{fixed timestep}, pembaruan selalu maju dengan potongan waktu yang
sama, misalnya $1/60$ detik, sedangkan sisa waktu yang belum terpakai
ditabung. Cara ini membuat simulasi dapat diulang persis dan perhitungan fisika
lebih stabil. Pada \istilah{variable timestep}, pembaruan memakai selang waktu
yang benar-benar terukur pada bingkai itu, sehingga panjang langkahnya
berubah-ubah. Program sesi ini termasuk golongan kedua: nilai dari
\code{clock.tick()} nanti dipakai langsung sebagai pengali kecepatan. Cara ini
sederhana dan memadai untuk permainan kecil, dengan satu catatan: bila
komputer tersendat dan satu bingkai menjadi sangat panjang, objek yang cepat
dapat melompati rintangan.

% ---------------------------------------------------------
\bagian{Tutorial}

Kerjakan langkah berikut berurutan. Setiap langkah adalah satu tindakan yang
selesai; periksa hasilnya sebelum melanjutkan.

\begin{langkah}
  \item \textbf{Siapkan folder sesi.} Salin folder \file{code/sesi-01} menjadi
        folder kerja baru, lalu buka folder itu di editor dan terminal. Dengan
        begitu Anda berangkat dari jendela yang sudah terbukti berjalan.

  \item \textbf{Tambahkan batas laju.} Buka \file{gudang_artefak/settings.py}
        dan tambahkan \code{FPS = 60} di dekat ukuran jendela. Angka 60 menjadi
        nilai bawaan yang dapat diubah lewat baris perintah nanti.

  \item \textbf{Tambahkan pengendali waktu.} Di dalam \code{Game.__init__},
        buat \code{self.clock = pygame.time.Clock()} dan
        \code{self.frames = 0}. Keduanya harus hidup selama permainan, jadi
        diletakkan sebagai atribut objek, bukan variabel di dalam
        perulangan.

  \item \textbf{Tulis siklus.} Ganti isi \code{run()} dengan perulangan
        \code{while self.running}. Di awal setiap putaran panggil
        \code{self.clock.tick(self.args.fps)}, lalu jalankan
        \code{handle_events()}, \code{draw()}, dan
        \code{pygame.display.flip()}. Tahap pembaruan belum berisi perintah
        karena keadaan permainan belum berubah; sisakan komentar penanda di
        antara \code{handle_events()} dan \code{draw()}. Karena perulangan
        berulang, pemanggilan \code{tick()} di awal putaran setara dengan
        menaruhnya di akhir putaran sebelumnya.

  \item \textbf{Tulis pembaca peristiwa.} Buat \code{handle_events()} yang
        mengulang \code{for event in pygame.event.get()}. Periksa
        \code{event.type == pygame.QUIT} dan tombol \code{K_ESCAPE}; keduanya
        mengubah \code{self.running} menjadi salah.

  \item \textbf{Tulis penayang.} Buat \code{draw()} yang mengisi latar dengan
        \code{self.screen.fill(S.BACKGROUND)}. Jangan menambahkan
        \code{flip()} di dalam \code{draw()}; penayangan tetap menjadi tugas
        \code{run()} agar urutan tahap mudah dibaca.

  \item \textbf{Tampilkan FPS pada judul.} Setelah \code{flip()}, tambahkan
        penghitung \code{self.frames}. Setiap 30 bingkai, panggil
        \code{pygame.display.set_caption()} dengan judul yang memuat
        \code{int(self.clock.get_fps())}. Membaca FPS setiap bingkai tidak
        perlu; jeda 30 bingkai membuat angkanya lebih tenang dibaca.

  \item \textbf{Jalankan program.} Dari folder sesi, ketik
        \code{python main.py}. Jendela harus tetap hidup, judulnya berubah
        menjadi \NamaPermainan\ diikuti angka FPS, dan menekan ESC atau tombol
        silang harus mengembalikan terminal dengan rapi.
\end{langkah}

\begin{ujibox}
Setelah Langkah 8, periksa empat hal berikut.
\begin{itemize}[leftmargin=1.4em, itemsep=0.2em]
  \item Judul jendela memuat angka FPS yang berkisar dekat 60, bukan 0 dan
        bukan angka yang tidak berubah sama sekali.
  \item Jendela tetap menanggapi ketika Anda mendiamkan tombol dan tetikus
        selama beberapa detik. Bila muncul tanda ``tidak menanggapi'',
        periksa pemanggilan \code{pygame.event.get()}.
  \item Menekan ESC menutup jendela, dan terminal menampilkan
        \code{Program selesai.}
  \item Latar tetap biru tua dan tidak meninggalkan jejak gambar lama. Bila
        ada jejak, periksa urutan \code{fill}, \code{draw}, lalu \code{flip}.
\end{itemize}
\end{ujibox}

Sekarang ubah batas laju dan amati akibatnya. Jalankan
\code{python main.py --fps 30}, lalu \code{python main.py --fps 120}. Bandingkan
angka pada judul jendela dengan angka yang Anda minta. Setelah itu jalankan
\code{python main.py --fps 0}: nilai nol berarti tanpa pembatasan, dan angka
FPS akan melonjak jauh di atas kemampuan layar Anda.

\begin{ujibox}
Sebelum melanjutkan, tunjukkan tiga bukti berikut kepada diri sendiri.
\begin{itemize}[leftmargin=1.4em, itemsep=0.2em]
  \item Pada \code{--fps 30}, angka pada judul berada di sekitar 30; pada
        \code{--fps 60}, di sekitar 60.
  \item Pada \code{--fps 0}, angka FPS jauh lebih tinggi dan pemakaian
        prosesor pada pengelola tugas meningkat tajam.
  \item Program dapat berhenti sendiri tanpa disentuh dengan
        \code{python main.py --frames 120} atau \code{--auto-close}; keduanya
        berguna untuk pengujian.
\end{itemize}
\end{ujibox}

% ---------------------------------------------------------
\bagian{Kode Sesi Ini}

Berkas berikut adalah isi checkpoint sesi 2. Bagian yang berubah dari sesi 1
adalah \file{settings.py} (penambahan \code{FPS}), \file{game.py} (siklus dan
pengendali waktu), dan \file{main.py} (pilihan \code{--fps} dan
\code{--frames}). Ketiklah sendiri di editor, lalu bandingkan hasilnya.

\tampilkodebaris[Python]{code/sesi-02/gudang_artefak/settings.py}{10}{13}
  {settings.py: judul, ukuran jendela, dan batas laju}

\tampilkodebaris[Python]{code/sesi-02/gudang_artefak/game.py}{19}{26}
  {game.py: clock dan penghitung bingkai disiapkan sekali}

\tampilkodebaris[Python]{code/sesi-02/gudang_artefak/game.py}{28}{58}
  {game.py: siklus permainan, pengendali waktu, dan judul FPS}

\tampilkodebaris[Python]{code/sesi-02/gudang_artefak/game.py}{60}{71}
  {game.py: pembacaan peristiwa dan penayangan latar}

\tampilkodebaris[Python]{code/sesi-02/main.py}{18}{30}
  {main.py: pilihan batas laju dan penghentian otomatis}

\checkpoint{02}{Program berjalan terus tanpa membeku, judul jendela
menampilkan laju bingkai terukur, dan program berhenti dengan rapi lewat
tombol silang, ESC, atau batas bingkai.}

% ---------------------------------------------------------
\bagian{Kesalahan Umum}

\begin{galatbox}
\begin{itemize}[leftmargin=1.4em, itemsep=0.35em]
  \item \textbf{Jendela membeku dan bertanda ``tidak menanggapi'.}
        Gejala: judul tidak berubah dan jendela tidak dapat ditutup dengan
        wajar. Periksa apakah \code{pygame.event.get()} dipanggil setiap
        bingkai; memakai \code{event.wait()} di dalam siklus menghasilkan
        gejala ini.
  \item \textbf{Angka FPS selalu 0.}
        Gejala: judul menampilkan \code{0 FPS} walaupun program berjalan.
        Periksa apakah \code{clock.tick()} benar-benar dipanggil dan apakah
        \code{get_fps()} dibaca setelah beberapa bingkai berjalan.
  \item \textbf{Prosesor bekerja penuh tanpa manfaat.}
        Gejala: kipas berbunyi keras dan pemakaian prosesor mendekati
        seratus persen. Periksa argumen \code{tick()}; memanggil
        \code{clock.tick()} tanpa angka tidak membatasi laju sama sekali.
  \item \textbf{Judul jendela tidak pernah berubah.}
        Gejala: judul tetap \NamaPermainan\ tanpa angka. Periksa blok
        yang memeriksa sisa bagi \code{self.frames} dengan 30, lalu pastikan
        \code{pygame.display.set_caption()} berada di dalam blok itu.
  \item \textbf{Layar berkedip atau isinya hitam.}
        Gejala: jendela terbuka tetapi tidak menampilkan latar yang benar.
        Periksa urutan \code{fill()}, \code{draw()}, lalu \code{flip()};
        menayangkan sebelum menggambar menghasilkan layar kosong.
  \item \textbf{Terminal tidak kembali setelah jendela ditutup.}
        Gejala: jendela hilang tetapi proses masih menggantung. Periksa
        penanganan \code{QUIT} dan pastikan perulangan benar-benar berakhir
        sebelum \code{pygame.quit()} dipanggil.
\end{itemize}
\end{galatbox}

% ---------------------------------------------------------
\bagian{Lembar Kerja}

Catat hasil pengamatan Anda pada tabel berikut. Isi kolom terakhir setelah
Anda membandingkan angka yang diminta dengan angka yang benar-benar tampil.

{\small
\renewcommand{\arraystretch}{1.4}
\begin{tabularx}{\linewidth}{@{}C{0.8cm} L{5.2cm} Y@{}}
\toprule
\textbf{No} & \textbf{Hasil yang dicatat} & \textbf{Keterangan Anda} \\
\midrule
1 & Batas laju bawaan dan angka FPS pada judul & \\[0.55cm]
2 & Angka FPS pada batas 30, 60, dan 120 & \\[0.9cm]
3 & Angka FPS dan pemakaian prosesor pada \code{--fps 0} & \\[0.9cm]
4 & Selisih antara batas yang diminta dan laju tercapai & \\[0.9cm]
5 & Gejala yang muncul ketika \code{event.get()} dihilangkan & \\[0.9cm]
6 & Satu kendala yang Anda hadapi dan cara menyelesaikannya & \\[0.9cm]
\bottomrule
\end{tabularx}
}

% ---------------------------------------------------------
\bagian{Checklist Sesi}

\begin{checklistbox}
\begin{checklist}
  \item Saya dapat menunjukkan jendela yang berjalan terus dan judulnya
        memuat angka FPS yang berubah.
  \item Saya dapat menjelaskan urutan tahap masukan, pembaruan, dan
        penayangan beserta alasan urutannya.
  \item Saya dapat menunjukkan bahwa program tetap menanggapi ketika pemain
        mendiamkan tombol dan tetikus.
  \item Saya dapat membatasi laju pada 60 FPS dan membandingkannya dengan
        30 FPS melalui judul jendela.
  \item Saya dapat menjelaskan dua pekerjaan \code{clock.tick()} dan mengapa
        perintah itu diletakkan pada awal atau akhir putaran siklus.
  \item Saya dapat menentukan apakah siklus yang saya tulis termasuk langkah
        tetap atau langkah berubah, beserta alasannya.
\end{checklist}
\end{checklistbox}

% ---------------------------------------------------------
\bagian{Pengait ke Sesi Berikutnya}

Siklus pada sesi ini sudah berjalan, tetapi layar masih menampilkan latar
polos. Tahap pembaruan belum mengubah apa pun. Pada sesi 3 Anda mengisi tahap
itu dengan seorang pemain berukuran 36 piksel yang bergerak mendatar pada
kecepatan 260 piksel per detik. Posisinya disimpan sebagai \code{Vector2}
bertipe pecahan, dan selang waktu yang telah Anda ukur pada sesi ini berubah
menjadi pengali kecepatan: $\vec{p}_{\text{baru}} = \vec{p} + \vec{v}\,\Delta
t$. Di sanalah angka FPS berhenti menjadi hiasan dan mulai menentukan jarak
yang ditempuh pemain setiap bingkai.
